% Template KLTN cho SV trường ĐHKHTN
% Liên hệ: nqminh@fit.hcmus.edu.vn
% Last update: 08/06/2016

% Chú ý: đọc các phần chú ý đóng khung của file này và chỉnh lại cho phù hợp.
% Trước khi build, xóa hết các file được tạo ra trong quá trình build trước đó, và build theo thứ tự: BIB > PDF > PDF.
% Nếu cập nhật tài liệu tham khảo, cũng cần build lại theo cách trên.

\documentclass[oneside,a4paper,14pt]{extreport}

% Font tiếng Việt
\usepackage[T5]{fontenc}
\usepackage[utf8]{inputenc}
\DeclareTextSymbolDefault{\DH}{T1}

% Tài liệu tham khảo
\usepackage[
	sorting=nty,
	backend=bibtex,
	defernumbers=true]{biblatex}
\usepackage[unicode]{hyperref} % Bookmark tiếng Việt
\addbibresource{References/references.bib}

\makeatletter
\def\blx@maxline{77}
\makeatother

% Chèn hình, các hình trong luận văn được để trong thư mục Images/
\usepackage{graphicx}
\graphicspath{ {Images/} }

% Chèn và định dạng mã nguồn
\usepackage{listings}
\usepackage{color}
\definecolor{codegreen}{rgb}{0,0.6,0}
\definecolor{codegray}{rgb}{0.5,0.5,0.5}
\definecolor{codepurple}{rgb}{0.58,0,0.82}
\definecolor{backcolour}{rgb}{0.95,0.95,0.92}
\lstdefinestyle{mystyle}{
    backgroundcolor=\color{backcolour},   
    commentstyle=\color{codegreen},
    keywordstyle=\color{magenta},
    numberstyle=\tiny\color{codegray},
    stringstyle=\color{codepurple},
    basicstyle=\footnotesize,
    breakatwhitespace=false,         
    breaklines=true,                 
    captionpos=b,                    
    keepspaces=true,                 
    numbers=left,                    
    numbersep=5pt,                  
    showspaces=false,                
    showstringspaces=false,
    showtabs=false,                  
    tabsize=2
}
\lstset{style=mystyle}

% Chèn và định dạng mã giả
\usepackage{amsmath}
\usepackage{algorithm}
\usepackage[noend]{algpseudocode}
\makeatletter
\def\BState{\State\hskip-\ALG@thistlm}
\makeatother

% Bảng biểu
\usepackage{multirow}
\usepackage{array}
\newcolumntype{L}[1]{>{\raggedright\let\newline\\\arraybackslash\hspace{0pt}}m{#1}}
\newcolumntype{C}[1]{>{\centering\let\newline\\\arraybackslash\hspace{0pt}}m{#1}}
\newcolumntype{R}[1]{>{\raggedleft\let\newline\\\arraybackslash\hspace{0pt}}m{#1}}

% Đổi tên mặc định
\renewcommand{\chaptername}{Chương}
\renewcommand{\figurename}{Hình}
\renewcommand{\tablename}{Bảng}
\renewcommand{\contentsname}{Mục lục}
\renewcommand{\listfigurename}{Danh sách hình}
\renewcommand{\listtablename}{Danh sách bảng}
\renewcommand{\appendixname}{Phụ lục}

% Dãn dòng 1.5
\usepackage{setspace}
\onehalfspacing

% Thụt vào đầu dòng
\usepackage{indentfirst}

% Canh lề
\usepackage[
  top=30mm,
  bottom=25mm,
  left=30mm,
  right=20mm,
  includefoot]{geometry}
  
% Trang bìa
\usepackage{tikz}
\usetikzlibrary{calc}
\newcommand\HRule{\rule{\textwidth}{1pt}}

% ========================================================================================= %
% CHÚ Ý: Thông tin chung về KLTN - sinh viên điền vào đây để tự động update các trang khác  %
% ========================================================================================= %
\newcommand{\tenSV}{Nguyễn~Văn~A~-~Trần~Văn~B} % Dấu ~ là khoảng trắng không được tách (các chữ nối với nhau bằng dấu ~ sẽ nằm cùng 1 dòng
\newcommand{\mssv}{1234567}
\newcommand{\tenKL}{Sử~dụng~LaTeX trong Khoá~luận~tốt~nghiệp} % Chú ý dấu ~ trong tên khóa luận
\newcommand{\tenGVHD}{Tên~Giáo~Viên}
\newcommand{\tenBM}{Công nghệ tri thức}

\begin{document}

\begin{titlepage}

\begin{center}
%ĐẠI HỌC QUỐC GIA THÀNH PHỐ HỒ CHÍ MINH\\
TRƯỜNG ĐẠI HỌC KHOA HỌC TỰ NHIÊN\\
\textbf{KHOA CÔNG NGHỆ THÔNG TIN}\\[2cm]


{ \Large \bfseries Hoàng Minh Thanh - Phan Minh Tâm\\[2cm] } 

%Tên đề tài Khóa luận tốt nghiệp/Đồ án tốt nghiệp

{ \Large \bfseries DỰ ĐOÁN LIÊN KẾT TRONG\\ĐỒ THỊ PHỨC\\KNOWLEDGE GRAPH EMBEDDING FOR LINKING PREDICTION\\[2cm] } 


%Chọn trong các dòng sau
\large KHÓA LUẬN TỐT NGHIỆP CỬ NHÂN\\
%\large ĐỒ ÁN TỐT NGHIỆP CỬ NHÂN\\
%\large THỰC TẬP TỐT NGHIỆP CỬ NHÂN\\
%Đưa vào dòng này nếu thuộc chương trình Chất lượng cao, hoặc lớp Cử nhân tài năng
\large CHƯƠNG TRÌNH HOÀN CHỈNH ĐẠI HỌC\\
%\large CHƯƠNG TRÌNH CHẤT LƯỢNG CAO\\
%\large CHƯƠNG TRÌNH CỬ NHÂN TÀI NĂNG\\[2cm]

\begin{tikzpicture}[remember picture, overlay]
  \draw[line width = 2pt] ($(current page.north west) + (2cm,-2cm)$) rectangle ($(current page.south east) + (-1.5cm,2cm)$);
\end{tikzpicture}

\vfill
Tp. Hồ Chí Minh, tháng 09/2020

\end{center}

\pagebreak



\begin{center}

TRƯỜNG ĐẠI HỌC KHOA HỌC TỰ NHIÊN\\
\textbf{KHOA CÔNG NGHỆ THÔNG TIN}\\[2cm]


{\large \bfseries Hoàng Minh Thanh - 18424062\\} 
{\large \bfseries Phan Minh Tâm - 18424059\\[2cm]}

%Tên đề tài Khóa luận tốt nghiệp/Đồ án tốt nghiệp

{ \Large \bfseries DỰ ĐOÁN LIÊN KẾT TRONG\\ĐỒ THỊ PHỨC\\KNOWLEDGE GRAPH EMBEDDING FOR LINKING PREDICTION\\[2cm] } 


%Chọn trong các dòng sau
\large KHÓA LUẬN TỐT NGHIỆP CỬ NHÂN\\
%\large ĐỒ ÁN TỐT NGHIỆP CỬ NHÂN\\
%Đưa vào dòng này nếu thuộc chương trình Chất lượng cao, hoặc lớp Cử nhân tài năng
\large CHƯƠNG TRÌNH HOÀN CHỈNH ĐẠI HỌC\\[2cm]
%\large CHƯƠNG TRÌNH CHẤT LƯỢNG CAO\\[2cm]
%\large CHƯƠNG TRÌNH CỬ NHÂN TÀI NĂNG\\[2cm]

\textbf{GIÁO VIÊN HƯỚNG DẪN}\\
Th.S Lê Ngọc Thành\\

\begin{tikzpicture}[remember picture, overlay]
  \draw[line width = 2pt] ($(current page.north west) + (2cm,-2cm)$) rectangle ($(current page.south east) + (-1.5cm,2cm)$);
\end{tikzpicture}

\vfill
Tp. Hồ Chí Minh, tháng 09/2020

\end{center}

\end{titlepage}
% Sasu trang Title, các bạn chèn nhận xét gủa GVHD và GVPB. Nhận xét sẽ được giáo vụ phát sau buổi bảo vệ để các bạn đóng quyển.

\pagenumbering{roman} % Đánh số i, ii, iii, ...

%\addcontentsline{toc}{chapter}{Lời cam đoan}
%\chapter*{Lời cam đoan}
\label{reassurances}

Tôi xin cam đoan đây là công trình nghiên cứu của riêng chúng tôi. Các số liệu và kết quả nghiên cứu trong luận văn này là trung thực và không trùng lặp với các đề tài khác.

\addcontentsline{toc}{chapter}{Lời cảm ơn}
\chapter*{Lời cảm ơn}
\label{thanks}

Tôi xin chân thành cảm ơn ...

\addcontentsline{toc}{chapter}{Đề cương chi tiết}
%Đây là template dùng cho đề cương đề tài tốt nghiệp
%Khoa Công nghệ Thông tin
%Trường Đại học Khoa học Tự nhiên, ĐHQG-HCM

%Liên hệ về mẫu LaTEX này: Thầy Bùi Huy Thông (bhthong@fit.hcmus.edu.vn)

\documentclass{article}[14pt]
\usepackage[utf8]{vietnam}
\usepackage{enumerate}
\usepackage{enumitem}
\usepackage{multicol}
\usepackage{listings}
\usepackage[left=2cm,right=2cm,top=2.5cm,bottom=2.5cm]{geometry}
\usepackage{verbatim}
\usepackage{graphicx}
\usepackage{url}
\usepackage{fancyhdr}
\usepackage{fancybox,framed}
\linespread{1.2}
\usepackage{lastpage}
\usepackage{floatrow}
\usepackage{floatrow}
\pagenumbering{arabic}
\usepackage{dirtytalk}
%\pagestyle{fancy}
\newfloatcommand{capbtabbox}{table}[][\FBwidth]

\usepackage{blindtext}
\usepackage{titlesec}
\usepackage[nottoc]{tocbibind}

\titleformat*{\section}{\LARGE\bfseries}
\titleformat*{\subsection}{\Large\bfseries}
\titleformat*{\subsubsection}{\large\bfseries}
%\addbibresource{ref.bib}


\begin{document}
    \begin{figure}[h]
        \begin{floatrow}
        \ffigbox{\includegraphics[scale = 0.1]{logo.png}}  
        {%
    
        }
        \capbtabbox{
            \begin{tabular}{l}
            \multicolumn{1}{c}{\textbf{\begin{tabular}[c]{@{}c@{}}TRƯỜNG ĐẠI HỌC KHOA HỌC TỰ NHIÊN\\KHOA CÔNG NGHỆ THÔNG TIN\end{tabular}}} \\ \\ \\
            \end{tabular}
        }
        {%
    
        }
        \end{floatrow}
    \end{figure}
    
    \begin{center}
        
        %Xác định loại đề tài tốt nghiệp tương ứng: Khóa luận, Thực tập, Đồ án
        \textbf{\Large ĐỀ CƯƠNG KHOÁ LUẬN TỐT NGHIỆP} \\ 
    \end{center}
    
    %\vspace{.5cm}
    
    \begin{center}
    %Tên đề tài phải VIẾT HOA
        
        \textbf{\huge Dự Đoán Liên Kết Trong Đồ Thị Tri Thức} 
        \\
        
    %Tên đề tài bằng tiếng Anh (nếu có)
    \vspace{.5cm}
        \textit{\textbf{\Large (Knowledge Graph Embedding for Link Prediction)}}
    \end{center}
    
    \vspace{.5cm}
    
    \Large
    \section{THÔNG TIN CHUNG}
    \begin{itemize}[label = {}]
        
        \item \textbf{Người hướng dẫn:} 
        %Thể hiện dạng: <Chức danh> <Họ và tên> (<Đơn vị công tác>)
        \begin{itemize}
            \item Ths. Lê Ngọc Thành (Khoa Công nghệ Thông tin)
        \end{itemize}{}
    
        
        \item \textbf{Nhóm Sinh viên thực hiện:}
        
        %Thể hiện dạng: <Họ và tên sinh viên> (MSSV: )
        \begin{enumerate}
        
            \item Phan Minh Tâm (18424059)
            \item Hoàng Minh Thanh (18424062)
        \end{enumerate}

       %Chọn loại thích hợp
        \item \textbf{Loại đề tài:} Nghiên cứu
        
        \item \textbf{Thời gian thực hiện:} Từ \textit{04/2020} đến \textit{10/2020}
        
        
    \end{itemize}
    
    \section{NỘI DUNG THỰC HIỆN}
    {

    %Mỗi mục dưới đây phải viết ít nhất là 5 câu mô tả/giới thiệu.
    
    \subsection{Giới thiệu về đề tài}
    
    Đồ thị tri thức (Knowledge Graphs-KG) là các biểu diễn cấu trúc của thông tin thế giới thực. Do khả năng mô hình hóa dữ liệu có cấu trúc, phức tạp theo cách máy tính có thể dễ dàng \say{hiểu được}, KG hiện đang được sử dụng rộng rãi trong nhiều lĩnh vực khác nhau, từ trả lời câu hỏi đến truy xuất thông tin và các hệ thống có thể suy luận dựa trên nội dung đã có.
    Việc phát triển một KG có thể được thực hiện bằng cách trích xuất các sự kiện mới từ các nguồn bên ngoài hoặc bằng cách suy ra các sự kiện còn thiếu từ những sự kiện đã có trong KG. Phương pháp tiếp cận, được gọi là Dự đoán liên kết (Link Prediction-LP).
    
    \subsection{Mục tiêu đề tài}
    %Phần này mô tả mục tiêu thực hiện đề tài.
    Cùng với nhiều kỹ thuật trí tuệ nhân tạo phát triển mạnh gần đây, đề tài tập chung nghiên cứu vào các khía cạnh của bài toán LP trên KG như đặc trưng tập dữ liệu, thuật toán, thực nghiệm đánnh giá các phương pháp cũng như các kỹ thuật khác nhau cùng tìm hiểu xem liệu những đặc trưng gì của tập dữ liệu hoặc các thuật toán khác nhau ảnh hưởng tới khả năng khái quát hóa của mô hình.
    \subsection{Phạm vi của đề tài}
    
    LP là một lĩnh vực nghiên cứu ngày càng sôi nổi gần đây đã phát triển mạnh mẽ từ sự bùng nổ của các kỹ thuật trong trí tuệ nhân tạo như: máy học (machine learning) và kỹ thuật học sâu (deep learning). Đề tài sẽ tập trung nghiên cứu các mô hình LP sử dụng KG làm nền tảng để tìm hiểu các biểu diễn dữ liệu với số chiều thấp còn được gọi là Knowledge Graph Embeddings, sau đó sử dụng chúng để suy ra các sự kiện, quan hệ mới.
    
    \subsection{Cách tiếp cận dự kiến}
    
    %Có thể bổ sung hình ảnh vào để làm rõ phương pháp hoặc cách tiếp cận dự kiến.
    Với mong muốn được tìm hiểu và nghiên cứu các phương pháp trí tuệ nhân tạo được áp dụng vào LG từ cổ điển đến hiện đại đề tài sẽ nghiên cứu một phương pháp cổ điển dựa trên luật (rule base \cite{burl}) và phương pháp còn lại dựa vào phương pháp học sâu (learning attention \cite{gat}). Đối với phương pháp cổ điển thuật toán cố gắng đưa các luật Horn \cite{wiki:Hornclause} sau đó khái quát hóa chúng thành các luật tổng quát trên tập dữ liệu huấn luyện. Phương pháp còn lại dựa vào kỹ thuật nhúng \cite{BordesUGWY13} đồ thị tri thức vào một không gian số chiều thấp dựa trên các láng giềng lân cận \cite{gat} sau đó học trên không gian này để đưa ra các dự đoán về các cạnh của đồ thị.
    
    \subsection{Kết quả dự kiến của đề tài}
    Tìm hiểu các phương pháp khác thực nghiệm và so sánh kết quả với phương pháp đang nghiên cứu với các phương pháp khác trên bốn tập dữ liệu được nghiên cứu nhiều nhất là FB15k \cite{toutanova2015observed}, FB15k-273 \cite{toutanova2015observed}, WN188 \cite{BordesUGWY13}, WN188RR \cite{GEmbedding}
    
    \subsection{Kế hoạch thực hiện}
    \begin{itemize}
        \item 4/2020 - 6/2020: Tìm hiểu các kiến thức liên quan về mạng nơ-ron,  KG. Đọc các  tài liệu bài báo liên quan tới đề tài.
        \item 6/2020 - 8/2020: Hiện thực các phương pháp, thuât toán đã đề xuất và tìm hiểu.
        \item 8/2020 - 9/2020: Tìm hiểu các phương pháp khác thực nghiệm và so sánh kết quả với phương pháp nghiên cứu.
        \item 9/2020 - 10/2020: Báo cáo và bảo vệ khóa luận tốt nghiệp.
    \end{itemize}
    }
    %TÀI LIỆU TRÍCH DẪN
    %Đây là ví dụ
    \bibliographystyle{ieeetr}
    \bibliography{sample}
    \nocite{*}

    \begin{table}[h]
    \centering
        \begin{tabular}{p{7cm}p{7cm}}
        \textbf{\begin{tabular}[c]{@{}c@{}}\\XÁC NHẬN\\CỦA NGƯỜI HƯỚNG DẪN\\ \textit{(Ký và ghi rõ họ tên)}\end{tabular}} & \textbf{\begin{tabular}[c]{@{}c@{}}\textit{TP. Hồ Chí Minh, ngày 20 tháng 09 năm 2020}\\NHÓM SINH VIÊN THỰC HIỆN\\\textit{(Ký và ghi rõ họ tên}) \end{tabular}}
        \end{tabular}
    \end{table}
    
\end{document}              

% Mục lục, danh sách hình, danh sách bảng
\addcontentsline{toc}{chapter}{Mục lục}
\tableofcontents
\listoffigures
\listoftables

\addcontentsline{toc}{chapter}{Tóm tắt}
\include{Appendix/tomtat}

\clearpage

\pagenumbering{arabic} % Đánh số 1, 2, 3, ...

% Các chương nội dung
%❖ Là chương mở rộng của tóm tắt
%❖ Nội dung bao gồm
%o Ngữ cảnh, vấn đề cần giải quyết là gì?
%o Vì sao vấn đề/bài toán quan trọng và thú vị?
%o Bài toán có gì khó? Vì sao cần phải giải quyết?
%o Những giải pháp (nghiên cứu, ứng dụng) đã giải quyết bài toán này?
%o Những giải pháp này có hạn chế, thiếu sót gì?
%o Giải pháp của bạn là gì? Kết quả thế nào?
%o Đóng góp của nghiên cứu/ứng dụng của bạn? 
%❖ Độ dài từ 3 đến 8 trang
%❖ Có thể chia ra các mục
%o Đặt vấn đề (Problem)
%o Mục tiêu (Objectives)
%o Giải pháp/Cách tiếp cận (An Approach)
%o Đóng góp (Contributions)
%o Bố cục (Outline)
%❖ Lưu ý, nội dung các chương sau phải khớp với mục tiêu, giải pháp, đóng góp ở chương này
\chapter{Giới thiệu}
\label{Chapter1}

Ngày nay đồ thị đã được ứng dụng vào mọi mặt của đời sống, với đồ thị về mạng xã hội (Facebook \cite{ugander2011anatomy}) thể hiện thông tin kết nối từng người với nhau, những nơi chúng ta đến, những thông tin chúng ta tương tác, đồ thị cũng được sử dụng làm cấu trúc trong hệ thống gợi ý video (Youtube \cite{baluja2008video}), trong mạng các chuyến bay, hệ thống định vị GPS, những tính toán khoa học hay thậm chí là kết nối não. Đồ thị tri thức của Google (Google's Knowledge Graph \cite{googlekg:2020}) được Google giới thiệu năm 2012 \cite{ji2020survey}, một loại đồ thị biểu diễn thông tin, là một trong những ứng dụng rõ ràng nhất về đồ thị tri thức cũng như cách dữ liệu được khai thác và biểu diễn trên đồ thị tri thức.

Khai thác đồ thị tri thức hiệu quả cung cấp cho người dùng hiểu sâu hơn về những gì đằng sau dữ liệu và từ đó có thể mang lại lợi ích cho nhiều ứng dụng trong thực tế. Tuy nhiên, trong thực tế, luôn có những tri thức mới được sinh ra mỗi ngày, thông tin thu được thường bị mất mát và không đầy đủ, từ đó nảy sinh ra vấn đề hoàn thiện đồ thị tri thức (knowledge graph completion) hay dự đoán liên kết (linking prediction) trong đồ thị tri thức.
Hầu hết các phương tiếp cận hiện nay là dự đoán một cạnh mới nối từ đỉnh này tới đỉnh khác. Với cách tiếp cận như vậy đồ thị tri thức có đầy đủ các tri thức nói cách khác làm cho đồ thị đày đặc nhờ tạo thêm cách cạnh nối. Nhưng như vậy chỉ mới giải quyết được vấn đề hoàn thành đồ thị, vấn đề thêm một (hoặc một lượng) tri thức mới vào đồ thị vẫn còn là một câu hỏi mở.

Hiện nay các bài toán liên quan đến hoàn thành đồ thị tri thức có hai cách tiếp cận chính là tối ưu hóa hàm mục tiêu tức là đưa ra dự đoán đựa trên có sai sót ít nhất như trong RuDiK\cite{ortona2018robust}, AMIE\cite{galarraga2015fast}, RuleN\cite{meilicke2018fine} liên quan tới các ứng dụng phân loại đỉnh, phân loại cạnh. Hoặc đưa ra một danh sách gồm \(k\) ứng viên với số điểm đại điện cho độ tin cậy giảm đần như trong các nghiên cứu TransE\cite{bordes2013translating}, ConvKB\cite{vu2019capsule} liên quan tới các hệ thống gợi ý (recommend system)...Cách tiếp cận của chúng tôi dựa trên cách tạo ra một danh sách gồm \(k\) ứng viên này.

Với mỗi cách tiếp cận trên có hai phương pháp chính đưa ra nghiên cứu một cách dựa trên luật như trong AnyBURL\cite{burl} hoặc dựa trên nhúng đồ thị như trong ConvE\cite{dettmers2017convolutional}, TransE\cite{bordes2013translating}, ComplEx \cite{trouillon2016complex}. Với mong muốn được tiếp cận các phương pháp có hệ thống nên chúng tôi chọn ở cả hai phương pháp để nghiên cứu thực hiện đề tài này. Đối với phương pháp dựa trên luật chúng tôi chọn phương pháp AnyBURL\cite{burl} còn với phương pháp dựa trên nhúng đồ thị chúng tôi chọn phương pháp KBAT \cite{nathani2019learning} sử dụng cơ chế chú  ý.

Đóng góp của chúng tôi trong phương pháp AnyBURL\cite{burl} gồm mã nguồn Python phương pháp AnyBURL. Cùng với đó chúng tôi cung cấp thêm hai chiến lược để thêm tri thức mới vào đồ thị mà chúng tôi gọi là online-to-offline là mở rộng của AnyBURL trong việc tạo ra các luật khi có một lượng (tập hợp) tri thức mới được thêm vào. Online-to-online sẽ tạo ngay các luật mới khi có một tri thức mới(cạnh) được thêm vào.

Với phương pháp dựa trên nhúng đồ thị chúng tôi sẽ trình bày lại về cơ chế chú ý (attention mechanisms \cite{vaswani2017attention}), cách cơ chế chú ý được áp dụng vào đồ thị tri thức bằng mô hình Mạng Đồ Thị Chú Ý (Graph Attention Network - GAT \cite{velivckovic2017graph}), mô hình KBAT \cite{nathani2019learning}, cũng như cải tiến của chúng tôi trên đồ thị dưa trên cải tiến mới nhất của cơ chế chú ý.
Đóng góp của chúng tôi trong phương pháp học sâu bao gồm mã nguồn trực tuyến trên Google Colab về cách cải tiến của mô hình chúng tôi, cũng như các phân tích của chúng tôi về các siêu tham số trong quá trình huấn luyện để đạt được kết quả tốt hơn.

% Hiện nay các bài toán liên quan đến hoàn thành đồ thị tri thức rất được quan tâm có bốn nhánh nghiên cứu chính như được nhắc đến trong nghiên cứu \cite{ampligraph} có hai nhánh nghiên cứu chính là nhúng đồ thị \cite{cai2018comprehensive} và dựa trên luật \cite{burl}. Với mong muốn được tiếp cận các phương pháp có hệ thống nên chúng tôi chọn ở mỗi nhánh một phương pháp làm chủ để chính cho các nghiên cứu và báo cáo này. Trong bài toán dự đự đoán liên kết cũng có hai phương pháp chính được đề xuất một là tối ưu hóa hàm mục tiêu đ


% Hiện nay các bài toán liên quan đến dự đoán liên kết đồ thị tri thức lớn rất được quan tâm có khoảng bốn nhánh nghiên cứu chính như được nhắc đến trong nghiên cứu \cite{ampligraph} một trong số đó là phương pháp dựa trên luật logic. Với mong muốn được tiếp cận các phương pháp từ đơn giản đến phức tạp nên chúng tôi chọn phương pháp này làm chủ để chính cho các báo cáo và nghiên cứu trong chương này. Với phương pháp này đưa ra một xếp hạng \(k\) ứng viên với một số điểm nhất định biểu thị cho sự chắc chán của dự đoán nó phù hợp với các hệ thống gợi ý(recommender system).

% Ngôn ngữ để viết và trình bày báo cáo khóa luận tốt nghiệp, đồ án tốt nghiệp, thực tập tốt nghiệp (sau đây gọi chung là báo cáo) là tiếng Việt hoặc tiếng Anh. 
% Trường hợp chọn ngôn ngữ tiếng Anh để viết và trình bày báo cáo,  sinh viên cần có đơn đề nghị, được cán bộ hướng dẫn (CBHD) đồng ý và nộp cho bộ phận Giáo vụ của Khoa vào thời điểm đăng ký đề tài để xin ý kiến.
% Báo cáo viết và trình bày bằng tiếng Anh phải có bản tóm tắt viết bằng tiếng Việt.


%Tóm tắt luận văn được trình bày nhiều nhất trong 24 trang in trên hai mặt giấy, cỡ chữ Times New Roman 11 của hệ soạn thảo Winword hoặc phần mềm soạn thảo Latex đối với các chuyên ngành thuộc ngành Toán.

%Mật độ chữ bình thường, không được nén hoặc kéo dãn khoảng cách giữa các chữ.
%Chế độ dãn dòng là Exactly 17pt.
%Lề trên, lề dưới, lề trái, lề phải đều là 1.5 cm.
%Các bảng biểu trình bày theo chiều ngang khổ giấy thì đầu bảng là lề trái của trang.
%Tóm tắt luận án phải phản ảnh trung thực kết cấu, bố cục và nội dung của luận án, phải ghi đầy đủ toàn văn kết luận của luận án.
%Mẫu trình bày trang bìa của tóm tắt luận văn (phụ lục 1).

\chapter{Trình bày báo cáo}
\label{Chapter2}

\section{Quy định chung}

Báo cáo phải được trình bày ngắn gọn, rõ ràng, mạch lạc, sạch sẽ, không được tẩy xóa, có đánh số trang, đánh số bảng biểu, hình vẽ, đồ thị. 

Nội dung báo cáo được phân thành các chương. Số thứ tự của các chương, mục được đánh số bằng hệ thống số Ả-rập, không dùng số La mã. Các mục và tiểu mục được đánh số bằng các nhóm hai hoặc ba chữ số, cách nhau một dấu chấm: số thứ nhất chỉ số chương, chỉ số thứ hai chỉ số mục, số thứ ba chỉ số tiểu mục.


%Báo cáo cần dùng LaTEX để viết và trình bày theo mẫu đã được cung cấp.

 Báo cáo trình bày sử dụng khổ giấy với việc canh lề như sau: Lề trên 3 cm, lề dưới 2,5 cm, lề trái 3 cm, lề phải 2 cm. Đánh số trang ở giữa bên dưới. Đánh số trang ở giữa bên dưới.

Font chữ dùng trong báo cáo (Times New Roman) với kích cỡ (size) 13-14pt, sử dụng chế độ dãn dòng (line spacing) chế độ 1.5 lines.

%Các bảng biểu trình bày theo chiều ngang khổ giấy thì đầu bảng là lề trái của trang. 


\section{Bố cục của báo cáo}

Nội dung của báo cáo tối thiểu 50 trang khổ A4 và không nên vượt quá 100 trang (không kể các trang bìa, lời cám ơn, mục lục, tài liệu tham khảo \ldots) theo trình tự như sau:

\begin{itemize}
\item MỞ ĐẦU (thường đặt tên là ``Giới thiệu''): Trình bày lý do chọn đề tài, mục đích, đối tượng và phạm vi nghiên cứu.
Mô tả bài toán mà đề tài giải quyết.
Bài toán này có gì hay?
Tại sao lại cần giải quyết bài toán này?
Bài toán này có gì khó?
Có những hướng nào để giải quyết bài toán này?
Những hướng giải quyết trước đây có những vấn đề gì chưa giải quyết được?
Các câu hỏi nghiên cứu mà đề tài trả lời hoặc những vấn đề mà đề tài sẽ giải quyết.
Các đóng góp của đề tài.

\item TỔNG QUAN (thường đặt tên là ``Các công trình liên quan''): Phân tích đánh giá các hướng nghiên cứu đã có của các tác giả trong và ngoài nước liên quan đến đề tài; nêu những vấn đề còn tồn tại (những vấn đề nào mà các công trình khác chưa giải quyết được); chỉ ra những vấn đề mà đề tài cần tập trung, nghiên cứu giải quyết.

\item NGHIÊN CỨU THỰC NGHIỆM HOẶC LÝ THUYẾT (thường đặt tên là ``Phương pháp đề xuất''): Trình bày cơ sở lý thuyết, lý luận, giả thiết khoa học và phương pháp nghiên cứu đã được sử dụng trong đề tài.

Nếu đề xuất hướng giải quyết mới, mô hình mới thì cần mô tả chi tiết cách giải quyết của mình (chi tiết tới mức người khác có thể dựa vào phần này mà cài đặt lại được đúng hoàn toàn phương pháp của mình đề ra).

\item TRÌNH BÀY, ĐÁNH GIÁ BÀN LUẬN VỀ CÁC KẾT QUẢ (thường đặt tên là ``Kết quả thí nghiệm''): Mô tả các kết quả nghiên cứu khoa học hoặc kết quả thực nghiệm.
Đối với  đề tài ứng dụng có kết quả là sản phẩm phần mềm phải có hồ sơ thiết kế, cài đặt,\ldots theo một trong các mô hình đã học (UML,\ldots).

Thông thường cần mô tả môi trường thí nghiệm trước như sử dụng dữ liệu nào, dùng độ đo nào để đánh giá, môi trường chạy thí nghiệm (cấu hình máy nếu cần phân tích thông tin về thời gian chạy thực nghiệm). Sau đó, nêu kết quả thực nghiệm, bàn luận và giải thích kết quả.

\item KẾT LUẬN VÀ HƯỚNG PHÁT TRIỂN (thường đặt tên là ``Kết luận''): Trình bày những kết quả đạt được, những đóng góp mới và những đề xuất mới, kiến nghị về những hướng nghiên cứu tiếp theo.

\item DANH MỤC TÀI LIỆU THAM KHẢO: Chỉ bao gồm các tài liệu được trích dẫn, sử dụng và đề cập tới để bàn luận trong báo cáo.
Phần này các bạn chuẩn bị 1 file BIB để lưu các tài liệu trích dẫn.
Khi các bạn trích dẫn một tài liệu nào đó, LaTeX sẽ tự động thêm vào danh mục tài liệu tham khảo giúp các bạn.
Các bạn xem hướng dẫn cách trích dẫn ở chương sau.

\item PHỤ LỤC: Phần này bao gồm nội dung cần thiết nhằm minh họa hoặc hỗ trợ cho nội dung báo cáo như số liệu, mẫu biểu, tranh ảnh,\ldots Phụ lục không được dày hơn phần chính của báo cáo.
Nếu có công trình công bố thì để vào phần phụ lục này.
\end{itemize}

\section{Bảng biểu, hình vẽ, phương trình}

%Những qui định dưới này các bạn có thể bỏ qua hoặc đọc để hiểu thêm.
%Những định dạng này LaTeX đều tự động giúp các bạn.
%Các bạn xem hướng dẫn chi tiết hơn ở chương sau.

Việc đánh số bảng biểu, hình vẽ, phương trình phải gắn với số chương; ví dụ hình 3.4 có nghĩa là hình thứ 4 trong Chương 3.
Mọi đồ thị, bảng biểu, hình vẽ lấy từ các nguồn khác phải được trích dẫn đầy đủ.

\subsection{Bảng biểu, hình vẽ}



Đầu đề của bảng biểu ghi phía trên bảng, đầu đề của hình vẽ ghi phía dưới hình.

Thông thường, những bảng ngắn và đồ thị phải đi liền với phần nội dung đề cập tới các bảng và đồ thị này ở lần thứ nhất.
Các bảng dài có thể để ở những trang riêng nhưng cũng phải tiếp theo ngay phần nội dung đề cập tới bảng này ở lần đầu tiên.
Các bảng rộng vẫn nên trình bày theo chiều đứng dài 297mm của trang giấy, chiều rộng của trang giấy có thể hơn 210mm.
Chú ý gấp trang giấy sao cho số và đầu đề của hình vẽ hoặc bảng vẫn có thể nhìn thấy ngay mà không cần mở rộng tờ giấy.
Tuy nhiên hạn chế sử dụng các bảng quá rộng này.

Đối với những trang giấy có chiều đứng hơn 297mm (bản đồ, bản vẽ,\ldots) thì có thể để trong một phong bì cứng đính bên trong bìa sau của báo cáo.

Các hình vẽ phải sạch sẽ bằng mực đen để có thể sao chụp lại; có đánh số và ghi đầy đủ đầu đề, cỡ chữ phải bằng cỡ chữ sử dụng trong báo cáo.

Khi đề cập đến các bảng biểu và hình vẽ phải nêu rõ số của hình và bảng biểu đó, ví dụ ``... được nêu trong Bảng 4.1'' hoặc ``xem Hình 3.2'' mà không được viết ``… được nêu trong bảng dưới đây'' hoặc ``trong đồ thị của X và Y sau''.

\subsection{Phương trình toán học}

Việc trình bày phương trình toán học trên một dòng đơn hoặc dòng kép tùy ý, tuy nhiên phải thống nhất trong toàn báo cáo.

Khi ký hiệu xuất hiện lần đầu tiên thì phải giải thích và đơn vị tính phải đi kèm ngay trong phương trình có ký hiệu đó.
Nếu cần thiết, danh mục của tất cả các ký hiệu, chữ viết tắt và nghĩa của chúng cần được liệt kê và để ở phần đầu của báo cáo.

Tất cả các phương trình cần được đánh số và để trong ngoặc đơn đặt bên phía lề phải.
Nếu một nhóm phương trình mang cùng một số thì những số này cũng được để trong ngoặc, hoặc mỗi phương trình trong nhóm phương trình (5.1) có thể được đánh số là (5.1.1), (5.1.2), (5.1.3).

\section{Viết tắt}

\textbf{Không lạm dụng việc viết tắt} trong báo cáo.
Chỉ viết tắt những từ, cụm từ hoặc thuật ngữ được sử dụng nhiều lần trong báo cáo.
Không viết tắt những cụm từ  dài, những mệnh đề; không viết tắt những cụm từ ít xuất hiện trong báo cáo.
Nếu cần viết tắt những từ thuật ngữ, tên các cơ quan, tổ chức,\ldots thì được viết tắt sau lần viết thứ nhất có kèm theo chữ viết tắt trong ngoặc đơn.
Nếu báo cáo có nhiều chữ viết tắt thì phải có bảng danh mục các chữ viết tắt (xếp theo thứ tự ABC) ở phần đầu báo cáo.

Nhắc lại: \textbf{không lạm dụng việc viết tắt} trong báo cáo.
Khi các bạn sử dụng từ viết tắt, người đọc sẽ phải lật lại những phần đã đọc, để tìm lại xem từ viết tắt đó nghĩa là gì.
Việc này sẽ làm chậm tốc độ đọc và sẽ khiến người đọc khó theo dõi báo cáo của bạn hơn.
Nếu có thể, hạn chế hoàn toàn việc dùng viết tắt.

\section{Tài liệu tham khảo và cách trích dẫn}

Mọi ý kiến, khái niệm có ý nghĩa, mang tính chất gợi ý không phải của riêng tác giả và mọi tham khảo khác phải được trích dẫn và chỉ ra nguồn trong danh mục Tài liệu tham khảo của báo cáo. Nguồn được trích dẫn phải được liệt kê chính xác trong danh mục Tài liệu tham khảo.

Việc trích dẫn, tham khảo chủ yếu nhằm thừa nhận nguồn của những ý tưởng có giá trị giúp người đọc theo được mạch suy nghĩ của tác giả, không làm trở ngại việc đọc.


Không trích dẫn những kiến thức phổ biến, mọi người đều biết cũng như không làm báo cáo nặng nề với những tham khảo trích dẫn.

Nếu không có điều kiện tiếp cận được một tài liệu gốc mà phải trích dẫn thông qua một tài liệu khác thì phải nêu ra trích dẫn này, đồng thời tài liệu gốc đó không được liệt kê trong danh mục tài liệu tham khảo của báo cáo.

\chapter{Hướng dẫn sử dụng template}
\label{Chapter3}

\section{Trích dẫn tài liệu}

Dùng lệnh $\textbackslash cite$ để trích dẫn một hoặc nhiều tài liệu tham khảo.
Tài liệu tham khảo có thể là trang web~\cite{Listings,HDLVThS}, bài báo khoa học~\cite{1994-Cavnar}, sách~\cite{1984-TeX-Knuth,2006-DDien,2006-NPTV}, bài tạp chí~\cite{1989-TED} hoặc các nguồn tham khảo khác. 
Lưu ý khi trích dẫn tài liệu tham khảo, cần viết câu sao cho bỏ phần trong cặp ngoặc vuông đi thì câu vẫn đầy đủ ý nghĩa.
Ví dụ, thay vì viết ``Nghiên cứu~\cite{1989-TED} chỉ ra rằng ... '' thì nên viết ``Nghiên cứu của Zhang~\cite{1989-TED} chỉ ra rằng ...''.
Một ví dụ khác, thay vì viết ``... như trong công trình nghiên cứu~\cite{1994-Cavnar}.'' thì nên viết ``... như trong công trình nghiên cứu của Cavnar và Trenkle~\cite{1994-Cavnar}.''

\section{Chèn mã nguồn}

Để chèn mã nguồn, cần dùng package listings~\cite{Listings}:

\begin{lstlisting}
\usepackage{listings}
\end{lstlisting}

Mã nguồn có thể được chèn trực tiếp như sau:

\begin{lstlisting}[language=Python]
print "Hello, World!"
\end{lstlisting}

hoặc chèn thông qua tập tin chứa mã nguồn trong thư mục $SourceCode$ như sau:

\lstinputlisting[language=C++]{SourceCode/hello.cpp}

Để chèn mã giả, cần dùng package algorithm~\cite{Algorithm}:

\begin{lstlisting}
\usepackage{algorithm}
\end{lstlisting}

Có thể chèn mã giả vào như sau:

\begin{algorithm}
\caption{My algorithm}\label{euclid}
\begin{algorithmic}[1]
\Procedure{MyProcedure}{}
\State $\textit{stringlen} \gets \text{length of }\textit{string}$
\State $i \gets \textit{patlen}$
\BState \emph{top}:
\If {$i > \textit{stringlen}$} \Return false
\EndIf
\State $j \gets \textit{patlen}$
\BState \emph{loop}:
\If {$\textit{string}(i) = \textit{path}(j)$}
\State $j \gets j-1$.
\State $i \gets i-1$.
\State \textbf{goto} \emph{loop}.
\State \textbf{close};
\EndIf
\State $i \gets i+\max(\textit{delta}_1(\textit{string}(i)),\textit{delta}_2(j))$.
\State \textbf{goto} \emph{top}.
\EndProcedure
\end{algorithmic}
\end{algorithm}

\section{Hình ảnh}

Để chèn hình ảnh, cần dùng package graphicx~\cite{Figures}:

\begin{lstlisting}
\usepackage{graphicx}
\end{lstlisting}

Hình \ref{fig:vd1}, hình \ref{fig:vd2} là một số ví dụ về chèn hình ảnh.

\begin{figure}[htp]
\centering
\includegraphics[width=6 cm]{images/logo-khtn.png}
\caption{Hình ví dụ 1}
\label{fig:vd1}
\end{figure}

\begin{figure*}[htp]
\centering
\includegraphics[width=40 mm]{images/logo-khtn.png}
\caption{Hình ví dụ 2}
\label{fig:vd2}
\end{figure*}

\section{Bảng biểu}

Để tạo bảng biểu, tham khảo thêm tại sharelatex.com~\cite{Tables}.
Bảng~\ref{tab:vd1} là một ví dụ về bảng.
Ngoài ra, có một số tool online~\footnote{Như trang http://www.tablesgenerator.com/} có thể được dùng để tạo bảng biểu một cách trực quan.

\begin{table}[ht]
\caption{Bảng ví dụ 1}

\label{tab:vd1}%
\begin{center}
\begin{tabular}{ |p{3cm}||p{3cm}|p{3cm}|p{3cm}|  }
 \hline
 \multicolumn{4}{|c|}{Country List} \\
 \hline
 Country Name     or Area Name& ISO ALPHA 2 Code &ISO ALPHA 3 Code&ISO numeric Code\\
 \hline
 Afghanistan   & AF    &AFG&   004\\
 Aland Islands&   AX  & ALA   &248\\
 Albania &AL & ALB&  008\\
 Algeria    &DZ & DZA&  012\\
 American Samoa&   AS  & ASM&016\\
 Andorra& AD  & AND   &020\\
 Angola& AO  & AGO&024\\
 \hline
\end{tabular}
\end{center}
\end{table}

\section{Công thức}

Công thức có thể chèn vào trong cùng một dòng như $ \sqrt{a^2+b^2} $ hoặc nằm trên dòng riêng như công thức~\ref{eu_eqn}.

\begin{equation} \label{eu_eqn}
x = a_0 + \frac{1}{a_1 + \frac{1}{a_2 + \frac{1}{a_3 + a_4}}}
\end{equation}

% Công trình của tác giả (nếu không có thì comment 02 dòng dưới)
\addcontentsline{toc}{chapter}{Danh mục công trình của tác giả}
\chapter*{Danh mục công trình của tác giả}
\label{Publish}

\begin{enumerate}
\item Tạp chí ABC
\item Tạp chí XYZ
\end{enumerate}

% In tài liệu tham khảo
\addcontentsline{toc}{chapter}{Tài liệu tham khảo}
\printbibheading[title={Tài liệu tham khảo}]

\printbibliography[heading=subbibliography, title={Tiếng Việt}, keyword=Viet, resetnumbers=true]

\DeclareNameAlias{sortname}{last-first}
\DeclareNameAlias{default}{last-first}

\printbibliography[heading=subbibliography, title={Tiếng Anh}, notkeyword=Viet, resetnumbers=4] 
% ===================================================================== %
% CHÚ Ý: phải gán lại resetnumbers=số tài liệu tham khảo tiếng Việt + 1 %
% ===================================================================== %

% Phần phụ lục
\appendix 

\chapter{Các siêu tham số tối ưu}
\label{Appendix1}

Trong phần này chúng tôi sẽ báo cáo lại tập hợp các siêu tham số tối ưu trên cả mô hình Chú ý và mô hình ConvKB.
Chúng tôi sử dụng tìm kiếm lưới theo Hits@10 để tìm các tham số tối ưu. Chúng tôi không chia nhỏ tập dữ liệu trong quá trình huấn luyện mô hình chú ý,
đối với mô hình dự đoán, chúng tôi sử dụng kích thước cho tất cả các tập dữ liệu. Cụ thể các siêu tham số chúng tôi sử dụng được trình bày trong bảng sau :

\begin{table}[htbp]
	\begin{center}
		\resizebox{\textwidth}{!}{%
	\begin{tabular}{llllllllll}
		\hline
		& $\mu$  & Weight decay & Epochs & negative ratio  & Dropouts  & $\alpha_{\text{LeakyRELU}}$ & $N_{\text{head}}$ & $D_{\text{final}}$ & $\gamma$ \\
		\hline
		FB15k     & 1e-3 & 1e-5 & 3000 & 2 & 0.3 & 0.2 & 2 & 200 & 1 \\
		FB15k-237 & 1e-3 & 1e-5 & 3000 & 2 & 0.3 & 0.2 & 2 & 200 & 1 \\
		WN18      & 1e-3 & 5e-6 & 3600 & 2 & 0.3 & 0.2 & 2 & 200 & 5 \\
		WN18RR    & 1e-3 & 5e-6 & 3600 & 2 & 0.3 & 0.2 & 2 & 200 & 5 \\
		\hline
	\end{tabular}
}
\end{center}
\end{table}
\chapter{Bảng thuật ngữ}
\label{Appendix2}

\printglossary[type=\acronymtype]

\end{document} 